%%%%%%%%%%%%%%%%%%%%%%%%
%
% $Autor: Sudhanshu Pukale $
% $Datum: 2025-03-24 $
% $Pfad: D:/Deutschland/Emden-leer/Ml-Attempt2/ML24-EX02-Pico4ML-main/ML24-EX02-Pico4ML-main/Pico4ML-BLE/Contents/General/RP2040.tex $
% $Version: 1 $
%
%%%%%%%%%%%%%%%%%%%%%%%%

\chapter{Microprocessor RP2040}
	
\section{Description}
	
	RP2040 is the debut microcontroller from Raspberry Pi. It brings our signature values of high performance, low cost, and ease of use to the microcontroller space.\cite{RaspberryPiRP2040:2024} With a large on-chip memory, symmetric dual-core processor complex, deterministic bus fabric, and rich peripheral set augmented with our unique Programmable I/O (PIO) subsystem, it provides professional users with unrivalled power and flexibility. With detailed documentation, a polished MicroPython port, and a UF2 bootloader in ROM, it has the lowest possible barrier to entry for beginner and hobbyist users. RP2040 is a stateless device, with support for cached execute-in-place from external QSPI memory.\cite{RaspberryPiRP2040:2024} This design decision allows you to choose the appropriate density of non-volatile storage for your application, and to benefit from the low pricing of commodity Flash parts. RP2040 is manufactured on a modern 40nm process node, delivering high performance, low dynamic power consumption, and low leakage, with a variety of low-power modes to support extended-duration operation on battery power.\cite{RaspberryPiRP2040:2024}




\begin{figure}
    \centering
    \begin{tikzpicture}
        
        % Define the chip size and colors
        \definecolor{chipblack}{RGB}{40, 40, 45}
        \definecolor{markings}{RGB}{180, 180, 180}
        
        % Main chip body - QFN package with 56 pins (7x7mm)
        \begin{scope}[xshift=0cm, yshift=0cm, rotate=0]
            % Bottom view chip
            \filldraw[chipblack] (-2,-2) rectangle (2,2);
            
            % Thermal pad at the center
            \filldraw[white] (-0.7,-0.7) rectangle (0.7,0.7);
            
            % Raspberry Pi logo
            \node at (-0.05, 0.25) {\includegraphics[width=0.4cm]{RP2040/RP2040.png}};
            
            % Text on the chip
            \node[text=markings, font=\sffamily\tiny] at (0, -0.2) {RP2-B2   21/25};
            \node[text=markings, font=\sffamily\tiny] at (-0.15, -0.4) {P4C754.00};
            
            % Pin markings (simplified representation of the 56 pins)
            % Top edge pins
            \foreach \x in {-1.7,-1.5,-1.3,-1.1,-0.9,-0.7,-0.5,-0.3,-0.1,0.1,0.3,0.5,0.7,0.9,1.1,1.3,1.5,1.7}
            \filldraw[markings] (\x,2) rectangle (\x+0.1,1.8);
            
            % Right edge pins
            \foreach \y in {-1.7,-1.5,-1.3,-1.1,-0.9,-0.7,-0.5,-0.3,-0.1,0.1,0.3,0.5,0.7,0.9,1.1,1.3,1.5,1.7}
            \filldraw[markings] (2,-\y) rectangle (1.8,-\y-0.1);
            
            % Bottom edge pins
            \foreach \x in {-1.7,-1.5,-1.3,-1.1,-0.9,-0.7,-0.5,-0.3,-0.1,0.1,0.3,0.5,0.7,0.9,1.1,1.3,1.5,1.7}
            \filldraw[markings] (\x,-2) rectangle (\x+0.1,-1.8);
            
            % Left edge pins
            \foreach \y in {-1.7,-1.5,-1.3,-1.1,-0.9,-0.7,-0.5,-0.3,-0.1,0.1,0.3,0.5,0.7,0.9,1.1,1.3,1.5,1.7}
            \filldraw[markings] (-2,-\y) rectangle (-1.8,-\y-0.1);
            
            % Orientation marking dot
            \filldraw[markings] (-1.6,1.6) circle (0.08);
        \end{scope}
        
    \end{tikzpicture}
    \caption{The RP2040 Schematic (adapted from \protect\cite{RaspberryPiRP2040:2024})}
    \label{fig:RP2040}
\end{figure}



\section{Hardware Interfaces}

\subsection{Communication Protocols}
\begin{itemize}
    \item \textbf{UART}: 2$\times$ Universal Asynchronous Receiver/Transmitter interfaces \cite{RaspberryPiRP2040:2024}
    \item \textbf{I$^2$C}: 2$\times$ Inter-Integrated Circuit interfaces supporting master and slave modes \cite{RaspberryPiRP2040:2024}
    \item \textbf{SPI}: 2$\times$ Serial Peripheral Interface buses \cite{RaspberryPiRP2040:2024}
    \item \textbf{USB}: USB 1.1 Host/Device controller with PHY \cite{RaspberryPiRP2040:2024}
\end{itemize}

\subsection{Input/Output}
\begin{itemize}
    \item \textbf{GPIO}: 30 General Purpose Input/Output pins \cite{RaspberryPiRP2040:2024}
    \item \textbf{ADC}: 4 analog inputs with 12-bit resolution (500 ksps) \cite{RaspberryPiRP2040:2024}
    \item \textbf{PWM}: 16 configurable PWM channels for motor control, servo positioning, and signal generation \cite{Schneider:2023}
\end{itemize}

\subsection{Advanced Interfaces}
\begin{itemize}
    \item \textbf{PIO}: 8 Programmable I/O state machines allowing implementation of custom digital interfaces \cite{RaspberryPiRP2040:2024}
    \item \textbf{Debug}: Serial Wire Debug (SWD) interface \cite{RaspberryPiRP2040:2024}
\end{itemize}

\section{Specifications}

\subsection{Processing Capabilities}
\begin{itemize}
    \item \textbf{CPU}: Dual-core ARM Cortex-M0+ processors \cite{RaspberryPiRP2040:2024}
    \item \textbf{Clock Speed}: Standard 133 MHz operation (certified for up to 200 MHz) \cite{RaspberryPiRP2040:2024}
    \item \textbf{Instruction Set}: ARM Thumb-2 instruction set \cite{RaspberryPiRP2040:2024}
\end{itemize}

\subsection{Memory}
\begin{itemize}
    \item \textbf{SRAM}: 264 KB on-chip SRAM divided into six banks \cite{Schneider:2023}
    \item \textbf{Flash}: External QSPI flash interface supporting up to 16 MB (typically 2 MB on Pico boards) \cite{Schneider:2023}
    \item \textbf{ROM}: Built-in boot ROM with support for UF2 bootloader \cite{Schneider:2023}
\end{itemize}

\subsection{Power Management}
\begin{itemize}
    \item \textbf{Voltage}: Core operating at 1.1V, I/O at 3.3V \cite{RaspberryPiRP2040:2024}
    \item \textbf{Power Modes}: Active, sleep, and dormant modes for energy efficiency \cite{RaspberryPiRP2040:2024}
    \item \textbf{Consumption}: Typical active power draw of 80mA at 133MHz \cite{RaspberryPiRP2040:2024}
\end{itemize}

\section{Constraints}

\subsection{Environmental}
\begin{itemize}
    \item \textbf{Temperature Range}: Industrial rating of -40\textdegree C to +85\textdegree C \cite{RaspberryPiRP2040:2024}
    \item \textbf{Humidity}: Non-condensing operation environment required \cite{RaspberryPiRP2040:2024}
\end{itemize}

\subsection{Performance}
\begin{itemize}
    \item \textbf{Real-time Capability}: Not certified for safety-critical applications but supports deterministic timing via PIO \cite{RaspberryPiRP2040:2024}
    \item \textbf{Interrupt Latency}: Typical interrupt response time of 12 clock cycles \cite{RaspberryPiRP2040:2024}
\end{itemize}

\subsection{Application}
\begin{itemize}
    \item \textbf{Certification}: Not rated for automotive or medical applications \cite{RaspberryPiRP2040:2024}
    \item \textbf{Reliability}: Designed for prototyping and commercial applications with appropriate environmental controls \cite{RaspberryPiRP2040:2024}
\end{itemize}

\section{Dimensions}

\begin{itemize}
    \item \textbf{Package}: QFN-56 package with 7$\times$7mm dimensions \cite{RaspberryPiRP2040:2024}
    \item \textbf{Pitch}: 0.4mm pin pitch \cite{RaspberryPiRP2040:2024}
    \item \textbf{Height}: 0.75mm package height \cite{RaspberryPiRP2040:2024}
\end{itemize}

\section{OS/Software Interfaces/Protocols}

\subsection{Development Environments}
\begin{itemize}
    \item \textbf{MicroPython}: Native support with simplified API for rapid development \cite{Schneider:2023}
    \item \textbf{C/C++ SDK}: Comprehensive SDK with peripheral libraries \cite{Schneider:2023}
    \item \textbf{Arduino}: Compatible with Arduino ecosystem through community-developed core \cite{Schneider:2023}
\end{itemize}

\subsection{Programming Interfaces}
\begin{itemize}
    \item \textbf{UF2 Bootloader}: Drag-and-drop programming via USB mass storage \cite{Schneider:2023}
    \item \textbf{SWD}: Hardware debugging through Serial Wire Debug protocol \cite{RaspberryPiRP2040:2024}
    \item \textbf{IDE Support}: Native support in Thonny, Visual Studio Code, and Arduino IDE \cite{Schneider:2023}
\end{itemize}

\subsection{System Software}
\begin{itemize}
    \item \textbf{RTOS Compatibility}: Compatible with FreeRTOS, RTX, and other embedded RTOS \cite{RaspberryPiRP2040:2024}
    \item \textbf{PIO Programming}: Custom assembly language for programmable I/O blocks \cite{RaspberryPiRP2040:2024}
\end{itemize}

\section{Data}

\subsection{Quality}
\begin{itemize}
    \item \textbf{ADC Resolution}: 12-bit resolution providing 4096 distinct levels \cite{RaspberryPiRP2040:2024}
    \item \textbf{Timing Precision}: Hardware timers with microsecond precision \cite{RaspberryPiRP2040:2024}
    \item \textbf{Jitter}: Low-jitter PWM generation suitable for motor control applications \cite{Schneider:2023}
\end{itemize}

\subsection{Quantity}
\begin{itemize}
    \item \textbf{Sampling Rate}: ADC capable of 500 thousand samples per second \cite{RaspberryPiRP2040:2024}
    \item \textbf{Data Throughput}: SPI interfaces supporting up to 33 MHz clock rates \cite{RaspberryPiRP2040:2024}
    \item \textbf{Buffer Sizes}: DMA controllers supporting configurable buffer transfers \cite{RaspberryPiRP2040:2024}
\end{itemize}

\subsection{Types}
\begin{itemize}
    \item \textbf{Analog}: Voltage input range of 0-3.3V on ADC pins \cite{RaspberryPiRP2040:2024}
    \item \textbf{Digital}: 3.3V CMOS logic levels with 5V-tolerant GPIO \cite{RaspberryPiRP2040:2024}
    \item \textbf{Serial}: UART supporting standard baud rates up to 921600 bps \cite{RaspberryPiRP2040:2024}
\end{itemize}

\subsection{Structure}
\begin{itemize}
    \item \textbf{Memory Mapping}: Memory-mapped peripheral registers \cite{RaspberryPiRP2040:2024}
    \item \textbf{DMA}: 12 channels of Direct Memory Access for efficient data movement \cite{RaspberryPiRP2040:2024}
    \item \textbf{Bus Architecture}: Advanced Peripheral Bus (APB) architecture for peripheral communication \cite{RaspberryPiRP2040:2024}
\end{itemize}

\section{Code with Simple Example}

The following MicroPython example demonstrates basic GPIO and ADC functionality on the RP2040:

\begin{lstlisting}[language=Python, caption={Basic RP2040 GPIO and ADC Example}]
    from machine import Pin, ADC
    import time
    
    # Configure LED on GPIO 25 (built-in on Pico)
    led = Pin(25, Pin.OUT)
    
    # Configure ADC on GPIO 26
    temp_sensor = ADC(26)
    
    # Configure a button input on GPIO 15 with pull-up
    button = Pin(15, Pin.IN, Pin.PULL_UP)
    
    def temperature_to_celsius(adc_value):
    # Convert ADC reading to voltage (0-3.3V)
    voltage = adc_value * 3.3 / 65535
    # Convert voltage to temperature (assuming LM35 sensor: 10mV/C)
    return voltage * 100
    
    while True:
    # Read temperature sensor
    raw_value = temp_sensor.read_u16()
    temperature = temperature_to_celsius(raw_value)
    
    # Print temperature every second
    print(f"Temperature: {temperature:.1f}C")
    
    # Check button state and toggle LED accordingly
    if button.value() == 0:  # Button pressed (active low with pull-up)
    led.toggle()         # Toggle LED state
    time.sleep(0.2)      # Debounce delay
    
    time.sleep(1)
\end{lstlisting}

This example demonstrates:
\begin{itemize}
    \item Digital output control with the onboard LED
    \item Analog input reading from a temperature sensor
    \item Digital input with pull-up resistor for button interfacing
    \item Basic timing control using the time module
\end{itemize}